\documentclass[10pt,a4paper]{amsart}
\usepackage[margin=19mm]{geometry}
\usepackage{amsmath,amssymb,mathtools}
\usepackage[scr=rsfs,cal=euler]{mathalfa}
\usepackage{xfrac}
\usepackage[unicode,hidelinks,bookmarks=false]{hyperref}
\newcommand{\ie}{i.e.\ }
\newcommand{\cf}{cf.\ }
\newcommand{\resp}{respectively\ }
\newcommand{\Subsection}{\subsection}
\providecommand{\nobreakdash}{\nobreak\hbox{-}}
\setlength{\emergencystretch}{2em}
\multlinegap0pt
\usepackage{amssymb}
\usepackage{csquotes}
\usepackage{mathtools}
\usepackage{tikz}
\usepackage{float}
\usepackage{xcolor}
\newtheorem{lemma}{Lemma}
\newtheorem{proposition}{Proposition}
\usepackage{cleveref}
\crefname{lemma}{Lemma}{Lemmas}
\crefname{proposition}{Proposition}{Propositions}
\newcommand{\asc}{\operatorname{asc}}
\newcommand{\des}{\operatorname{des}}
\newcommand{\is}{\mathsf{is}}
\newcommand{\ls}{\mathsf{ls}}
\newcommand{\s}{\mathfrak{S}}
\title{\large\textbf{Descent-restricted subsequences via RSK and evacuation}}
\author{Krishna Menon, Anurag Singh}
\date{Source: arXiv:2602.04767v2 (CC BY 4.0)}
\begin{document}
\maketitle

\noindent\emph{Source and licence.} \large\textbf{Descent-restricted subsequences via RSK and evacuation}. Version arXiv:2602.04767v2 (CC BY 4.0), distributed by arXiv under the Creative Commons Attribution 4.0 International licence, http://creativecommons.org/licenses/by/4.0/, as recorded in the arXiv licence metadata for this exact version and shown on the abstract page https://arxiv.org/abs/2602.04767v2. Full licence notice: \texttt{licenses/arxiv-2602.04767-CC-BY-4.0.txt}.\par\smallskip

\noindent\textit{Editorial scope.} Counted unit: the article's proposition gs=isplus1 (source0.tex lines 550--552). The article's complete original proof of it is delivered with it (source0.tex lines 554--580, one \texttt{proof} environment of 27 lines). No mathematics is written, repaired or re-derived: the statement and the proof are the article's own lines, in the article's own notation, and no retained line is substituted or re-flowed.  Retained uncounted as the source-local context the proof needs: lines 507--510; lines 527--532.  Nothing was omitted from any retained span, so the package carries no omission record.  No retained line contains a citation, so the wrapper carries no bibliography and no bibliography entry.  No retained line contains a figure environment, an image-inclusion command or drawing code, so no image is delivered and the package declares no asset dependency.  Editorial formatting only, in addition: an AMS \texttt{amsart} preamble that restates the article's own declarations and macros used by the retained lines, namely the environments \texttt{lemma}, \texttt{proposition} on separate counters, together with the standard packages those lines need.  The retained environments print in this document as Lemma 1, Proposition 1 in order of appearance, whereas the article prints the same items with the numbering of its own full document.  The complete native source of the article as distributed in the arXiv e-print for this exact version is bundled unchanged as \texttt{sources/193994/source0.tex}.
\begin{lemma}\label{lsdnonzero}
	For any permutation $\pi$, we have $\ls_d(\pi) \geq d + 1$ for all
	$d \leq \des(\pi)$ and $\ls_d(\pi) = 0$ for all $d > \des(\pi)$. Moreover, for any $d_1 < d_2 \leq \des(\pi)$, we have $\ls_{d_1}(\pi) < \ls_{d_2}(\pi)$.
\end{lemma}

\begin{lemma}\label{lsdbound}
    For any permutation $\pi$ and $d \leq \des(\pi)$, we have
    \begin{equation*}
        \is(\pi) + d \leq \ls_d(\pi) \leq \operatorname{min}\{\is_{d + 1}(\pi), \asc(\pi) + d + 1\}.
    \end{equation*}
\end{lemma}

\begin{proposition}\label{gs=isplus1}
    For any permutation $\pi$ and $1 \leq d \leq \des(\pi)$, we have $\ls_d(\pi) = \is(\pi) + d$ if and only if $\asc(\pi) = \is(\pi) - 1$.
\end{proposition}

\begin{proof}
    If $\asc(\pi) = \is(\pi) - 1$, then \Cref{lsdbound} immediately gives $\ls_d(\pi) = \is(\pi) + d$. 
    Let $\pi \in \s_n$. 
    Note that if there exists $1 \leq i < j \leq n$ such that $\pi(i) < \pi(j)$, then there exists an ascent at some index in $[i, j - 1]$. 
    Using an increasing subsequence of length $\is(\pi)$, this shows that $\asc(\pi) \geq \is(\pi) - 1$.
    
    Hence, by \Cref{lsdbound}, we have to show that if $\asc(\pi) \geq \is(\pi)$, then $\ls_d(\pi) \geq \is(\pi) + d + 1$. 
    We only have to prove this for the case $d = 1$. 
    Since if $\ls_1(\pi) \geq \is(\pi) + 2$, from \Cref{lsdnonzero}, we get $\ls_1(\pi) < \ls_2(\pi) < \cdots < \ls_d(\pi)$, which gives us $\ls_d(\pi) \geq \is(\pi) + d + 1$, as required.

    Set $k = \is(\pi)$ and suppose that $\asc(\pi) \geq k$. 
    We will show that $\ls_1(\pi) \geq k + 2$ by constructing a subsequence of length $k + 2$ with exactly one descent. 
    Suppose that $I = \{i_1 < i_2 < \cdots < i_k\}$ is such that $\pi_I$ is an increasing subsequence of $\pi$.

    If $\pi$ has an ascent $j \in [1, i_1 - 1] \cup [i_k, n]$, then $\pi_{I \cup \{j, j + 1\}}$ is a subsequence of length $k + 2$ with one descent. 
    Note that such a $j$ can never be equal to $i_1 - 1$ or $i_k$ since this would contradict the fact that $\is(\pi) = k$.

    If $\pi$ has no ascent of the form mentioned above, then since $\asc(\pi) \geq k$, there must exist some $m \in [k - 1]$ such that $\pi$ has at least two ascents in $[i_m, i_{m + 1} - 1]$. 
    Note that $\is(\pi) = k$ implies that for any $i \in [i_m, i_{m + 1} - 1]$, $\pi(i) < \pi(i_m)$ or $\pi(i) > \pi(i_{m + 1})$. 
    The fact that there are at least two ascents in $[i_m, i_{m + 1} - 1]$ means that there exist $i, j \in [i_m, i_{m + 1} - 1]$ where $i < j$ such that
    \begin{enumerate}
        \item $\pi(i) > \pi(i_{m + 1})$ and $\pi(j) < \pi(i_m)$, or
        \item $\pi(i_{m + 1}) < \pi(i)< \pi(j)$, or
        \item $\pi(i) < \pi(j) < \pi(i_m)$.
    \end{enumerate}
    In each case, $\pi_{I \cup \{i, j\}}$ is a subsequence of length $k + 2$ with one descent.
\end{proof}

\end{document}
